\documentclass[12pt,a4paper]{article}

\usepackage[margin=1in]{geometry}
\usepackage{setspace}
\usepackage{longtable}
\usepackage{booktabs}
\usepackage{tabularx}
\usepackage{array}
\usepackage{needspace}
\usepackage{hyperref}
\usepackage[T1]{fontenc}
\usepackage[utf8]{inputenc}

\setstretch{1.12}
\setlength{\parskip}{6pt}
\setlength{\parindent}{0pt}

\begin{document}

\begin{center}
    {\LARGE \textbf{Backend Testing and Verification Report}}\\[0.4cm]
    {\large Resource Allocation System}\\[0.3cm]
    Najib Attar\\
    May 2026
\end{center}

\section{Introduction}

This report summarizes the backend testing and verification phase for the C++ Resource Allocation System. The system supports users, spaces, requests, allocations, request metadata, lifecycle history, rule-based validation, strategy-based allocation, CSV input/output, and configuration-based strategy selection.

The old \texttt{doctest}-based tests were no longer compatible with the current architecture because the source code had changed significantly. In particular, the user model became polymorphic, \texttt{User} became abstract, rule classes moved to an \texttt{evaluate(...)} interface, and request constructors were extended with metadata fields. Therefore, the outdated tests were replaced with a lightweight local test harness under the \texttt{tests/} directory.

The purpose of this testing phase was not only to show that normal examples work, but also to test defect-prone cases, regression risks, malformed input, strategy behavior, CSV loading/exporting, and backend execution.

\section{Testing Objective}

The main objective was to perform defect-oriented backend testing after many architectural changes. The tests were designed to verify both small unit behavior and larger integration behavior.

The testing goals were:

\begin{itemize}
    \item verify current backend behavior against the actual source code;
    \item detect regressions caused by architectural refactoring;
    \item test invalid and malformed data, not only happy paths;
    \item verify strategy ordering behavior for Greedy and Priority strategies;
    \item confirm CSV input and output behavior;
    \item confirm that backend execution still produces allocation and result files;
    \item identify remaining risks without claiming full mathematical correctness.
\end{itemize}

\section{Test Environment}

The tests were executed in a Windows PowerShell environment using \texttt{g++} with C++17.

\begin{center}
\begin{tabularx}{\textwidth}{@{} p{4cm} X @{}}
\toprule
\textbf{Item} & \textbf{Value} \\
\midrule
Operating environment & Windows PowerShell \\
Compiler & \texttt{g++} \\
Language standard & \texttt{-std=c++17} \\
Test runner & Local lightweight test harness in \texttt{tests/} \\
External unit test dependency & None \\
Main test command & \texttt{powershell -ExecutionPolicy Bypass -File tests\textbackslash run\_tests.ps1} \\
\bottomrule
\end{tabularx}
\end{center}

The local test framework is implemented in \texttt{tests/test\_framework.h}. It provides simple test registration, assertions, pass/fail output, and a summary of total passed and failed tests.

\section{Testing Strategy}

The testing strategy combines unit testing, integration testing, backend smoke testing, negative testing, and regression testing.

\begin{center}
\begin{tabularx}{\textwidth}{@{} p{4cm} X @{}}
\toprule
\textbf{Test type} & \textbf{Purpose} \\
\midrule
Unit tests & Verify isolated behavior of models, factories, rules, and services. \\
Integration tests & Verify interaction between CSV loading, request processing, allocation, and CSV export. \\
Backend smoke test & Compile and run the real backend executable and verify generated output files. \\
Negative tests & Use malformed, missing, invalid, and conflicting data to try to break the system. \\
Regression tests & Protect behavior introduced by previous iterations, such as metadata, lifecycle history, strategy selection, and invalid request handling. \\
Compiler warning check & Compile with stronger warning flags to detect suspicious code. \\
\bottomrule
\end{tabularx}
\end{center}

\section{Verified User Roles and Priorities}

The current source code defines the following user roles and priorities:

\begin{center}
\begin{tabularx}{0.65\textwidth}{@{} X c @{}}
\toprule
\textbf{Role} & \textbf{Priority} \\
\midrule
Student & 1 \\
TeachingAssistant & 2 \\
Staff & 3 \\
Instructor & 4 \\
Administrator & 5 \\
\bottomrule
\end{tabularx}
\end{center}

The \texttt{TeachingAssistant} role is implemented in the source tree. It is supported by \texttt{UserFactory}, has its own priority and permissions, and appears in \texttt{data/users.csv}. This confirms that \texttt{TeachingAssistant} is not only a test assumption, but part of the actual system model.

\section{Unit Testing}

The unit tests verify individual backend components. These include \texttt{TimeSlot}, user subclasses, \texttt{UserFactory}, \texttt{RequestFactory}, rules, availability logic, and strategy factory behavior.

Examples of unit-tested behavior include:

\begin{itemize}
    \item \texttt{TimeSlot} detects overlap on the same day;
    \item boundary intervals such as 10:00--12:00 and 12:00--14:00 do not overlap;
    \item invalid time ranges are rejected;
    \item \texttt{UserFactory} creates supported user subclasses correctly;
    \item unknown user roles are handled safely;
    \item role priorities and permissions match the current implementation;
    \item \texttt{RequestFactory} creates \texttt{OneTimeRequest}, \texttt{RecurringRequest}, \texttt{ExamRequest}, and \texttt{CommitteeMeetingRequest};
    \item malformed request time data creates \texttt{InvalidRequest};
    \item request title, purpose, and lifecycle history are preserved;
    \item capacity, feature, status, location, role, and availability rules behave correctly.
\end{itemize}

\section{Integration Testing}

Integration tests use temporary CSV files rather than relying only on the normal \texttt{data/} folder. This makes the tests reproducible and isolates them from manual demo data.

The integration tests verify that:

\begin{itemize}
    \item users, spaces, and requests can be loaded through \texttt{DataController};
    \item loaded requests can be processed through \texttt{AllocationService};
    \item approved and rejected statuses are assigned correctly;
    \item \texttt{allocations.csv} can be exported;
    \item \texttt{request\_results.csv} can be exported;
    \item exported results include title, purpose, requester role, priority, status, rejection reason, and lifecycle history;
    \item invalid rows are represented as \texttt{InvalidRequest} objects rather than causing crashes.
\end{itemize}

\section{System / Backend Smoke Testing}

The PowerShell runner also performs a backend smoke test before running the unit and integration test suite.

The smoke test performs the following steps:

\begin{enumerate}
    \item compile all backend source files with \texttt{g++ -std=c++17};
    \item run the generated backend executable;
    \item check the backend exit code;
    \item verify that \texttt{data/allocations.csv} exists and is not empty;
    \item verify that \texttt{data/request\_results.csv} exists and is not empty.
\end{enumerate}

This confirms that the real executable still runs successfully using the current normal data files.

\section{Negative and Regression Testing}

The test suite includes defect-oriented cases designed to break hidden assumptions.

The added negative and regression cases include:

\begin{itemize}
    \item missing users, spaces, and requests files;
    \item empty users, spaces, and requests files;
    \item unknown user roles;
    \item unknown space types;
    \item duplicate request IDs becoming \texttt{InvalidRequest};
    \item malformed request rows becoming \texttt{InvalidRequest};
    \item extra request CSV columns being tolerated;
    \item zero-length time slots being rejected;
    \item negative and reversed time data being rejected;
    \item exact capacity passing;
    \item one-over capacity failing;
    \item empty feature and building preferences passing;
    \item one-time external availability conflicts failing;
    \item recurring external availability conflicts failing;
    \item recurring internal self-conflicts failing;
    \item empty configuration values falling back to the Greedy strategy.
\end{itemize}

\section{Test Coverage Matrix}

\begin{longtable}{@{} p{5.2cm} p{2.2cm} p{7.0cm} @{}}
\toprule
\textbf{Feature / Module} & \textbf{Covered} & \textbf{Evidence} \\
\midrule
\endfirsthead

\toprule
\textbf{Feature / Module} & \textbf{Covered} & \textbf{Evidence} \\
\midrule
\endhead

TimeSlot overlap, boundaries, invalid ranges & Yes & Same-day overlap, different-day non-overlap, half-open boundary behavior, invalid range rejection. \\
UserFactory and polymorphic users & Yes & Student, TeachingAssistant, Staff, Instructor, and Administrator creation verified. \\
TeachingAssistant support & Yes & Class, factory support, role name, priority, permissions, and data presence verified. \\
User priorities and permissions & Yes & Role priorities and allowed/requested space types are tested. \\
RequestFactory valid/invalid creation & Yes & Valid request subtypes and invalid malformed requests are tested. \\
Request metadata and history & Yes & Title, purpose, default metadata, approval, rejection, and lifecycle events are tested. \\
DataLoader missing/empty/malformed CSVs & Yes & Missing files, empty files, malformed rows, unknown roles, unknown spaces, and duplicate user/space IDs are tested. \\
DataController temporary CSV integration & Yes & Temporary users, spaces, requests, auxiliary busy-slot paths, auxiliary participant paths, processing, and exports are tested. \\
InvalidRequest behavior & Yes & Unknown request type, malformed row, invalid user, invalid space, and bad time data are tested. \\
Capacity, Feature, Status, Location, UserRole rules & Yes & Pass and fail behavior for each rule category is tested. \\
AvailabilityRule conflicts & Yes & One-time and recurring external conflicts are tested. \\
RuleEngineFacade through AllocationService & Yes & Requests are processed through AllocationService, which uses RuleEngineFacade internally. \\
Greedy strategy order & Yes & Conflicting requests verify input-order processing. \\
Priority strategy order and stability & Yes & Descending priority and same-priority stable order are tested. \\
AllocationStrategyFactory fallback & Yes & Greedy, Priority, uppercase input, empty input, and unknown strategy fallback are tested. \\
Config loading fallback & Yes & Greedy, Priority, missing key, invalid value, empty value, and missing config are tested. \\
CSV exports & Yes & Allocations and request results are exported and checked for expected content. \\
Backend smoke execution & Yes & Real backend executable is compiled and run before the test suite. \\
\bottomrule
\end{longtable}

\section{Test Results}

The main test command was:

\begin{verbatim}
powershell -ExecutionPolicy Bypass -File tests\run_tests.ps1
\end{verbatim}

The result was:

\begin{verbatim}
Running 56 test(s)
Result: 56 passed, 0 failed
\end{verbatim}

The runner also compiled and ran the backend smoke executable before running the unit and integration tests.

\section{Compiler Warning Check}

A compiler warning check was performed using:

\begin{verbatim}
g++ -std=c++17 -Wall -Wextra
\end{verbatim}

No compiler warnings were reported during this check.

\section{AddressSanitizer Check}

An AddressSanitizer build was attempted using:

\begin{verbatim}
-fsanitize=address
\end{verbatim}

However, AddressSanitizer was not available in the current MinGW environment because the linker could not find \texttt{-lasan}. Therefore, memory safety is not tool-verified in this testing phase.

\section{Defects and Risks Found}

No allocation-behavior defect requiring an immediate redesign was found during this pass.

The testing phase originally identified several maintenance risks:

\begin{itemize}
    \item \texttt{DataLoader} allowed duplicate user IDs and duplicate space IDs. This did not crash the system, but it could create ambiguous lookup behavior.
    \item \texttt{DataController::loadAllData(...)} accepted custom users, spaces, and requests paths, but still hardcoded auxiliary paths such as user busy slots and request participants.
    \item \texttt{TeachingAssistant} was implemented in code and appeared in data, but some documentation did not consistently describe this role.
\end{itemize}

After the report, a small data-loading robustness cleanup addressed the first two risks:

\begin{itemize}
    \item duplicate user IDs are now skipped with a clear warning while preserving the first valid user;
    \item duplicate space IDs are now skipped with a clear warning while preserving the first valid space;
    \item \texttt{DataController::loadAllData(...)} now accepts configurable auxiliary paths for user busy slots and request participants while preserving default paths for normal backend execution.
\end{itemize}

The remaining documentation risk is addressed by updating the project README to reflect the current roles, request types, rules, strategies, committee meeting support, and testing location.

\section{Remaining Limitations}

The current backend test suite is stronger than before, but it is not a proof of full correctness.

Remaining limitations include:

\begin{itemize}
    \item no AddressSanitizer or memory-leak tooling is available in the current environment;
    \item no fuzz testing is performed;
    \item no exhaustive mathematical verification of allocation algorithms is performed;
    \item future multi-room exam allocation edge cases may require additional specialized tests;
    \item generated CSV files are verified for existence and key content, but not every column is exhaustively validated in every scenario.
\end{itemize}

\section{Summary}

This backend testing phase replaced outdated tests with a working local test harness and expanded the suite into a defect-oriented regression test set. The current suite covers core models, factories, rules, strategies, configuration loading, CSV loading, CSV exporting, invalid request behavior, request metadata, lifecycle history, and backend smoke execution.

The final test result was:

\begin{center}
\textbf{56 tests passed, 0 tests failed.}
\end{center}

No production source-code defect requiring an allocation-behavior redesign was found. The backend is ready for formal testing reporting and demonstration, while still having clearly identified limitations related to unavailable memory-safety tooling, lack of fuzz testing, and non-exhaustive mathematical verification.

\end{document}
